%  02_related_work.tex
\section{Related Work}
\label{sec:related}

\subsection{Animal re-identification}

Automated individual identification from imagery has been explored for a range of
species~\cite{schneider2022reidentification}.
Early work on sea turtles, whale sharks, and manta rays~\cite{ravela2004visual,
moskvyak2021robust} used hand-crafted texture descriptors.  More recently, deep
convolutional networks have been adapted from human face recognition to animal identification,
typically via fine-tuning on species-specific datasets.  A persistent gap in many published
systems is the use of identity-leaky evaluation splits, which makes cross-paper comparisons
unreliable and overstates deployment-ready performance.  Our work explicitly addresses
this gap.

\subsection{Metric learning and face verification}

The Siamese network, originally proposed for signature verification~\cite{bromley1994signature}
and popularised for one-shot classification~\cite{koch2015siamese}, learns an embedding
space where pairs of the same identity are pulled together and different-identity pairs
are pushed apart.  The similarity score is computed directly from the pair, which aligns
the training objective with the verification task.
FaceNet~\cite{schroff2015facenet} extended this idea with triplet mining on large-scale
human face datasets, achieving strong open-set performance.

\subsection{Angular margin softmax losses}

A parallel line of work treats face recognition as a closed-set classification problem
at training time and uses the resulting embedding space for open-set retrieval at test time.
SphereFace~\cite{liu2017sphereface} introduced a multiplicative angular margin;
CosFace~\cite{wang2018cosface} replaced it with an additive cosine margin;
ArcFace~\cite{deng2019arcface} used an additive angular margin directly in the angle space,
giving sharper decision boundaries and better calibration.
These losses share a common structure: the logit for the ground-truth class is modified
to require a larger cosine distance from the nearest other class, promoting compact,
well-separated clusters in the hyperspherical embedding space.

Sub-center ArcFace~\cite{sun2020subcenter} relaxes the assumption of a single clean
prototype per class, allowing multiple sub-centres to handle noisy labels.
AdaFace~\cite{kim2022adaface} further adapts the margin magnitude to the estimated quality
of each image, using the \emph{feature norm} as a proxy: low-quality images (whose
embeddings have smaller norms) receive a smaller margin, reducing their influence on
the decision boundary.

While these angular-margin methods dominate public human face benchmarks (where training
sets contain millions of identities), their behaviour under very few training identities---as
in ecological datasets---has not been systematically studied.  Our experiments fill this gap.

\subsection{Biometric evaluation metrics}

Standard face recognition benchmarks report rank-1 identification accuracy and
verification equal-error rate (EER).  For security-critical biometric applications
the operating metric is \emph{TAR at a fixed FAR} (true acceptance rate at false acceptance
rate), reflecting real-world deployment constraints.
In our setting, where test partitions contain at most a few hundred image pairs, FAR
targets such as $10^{-3}$ are effectively degenerate (budgeting fewer than one false match),
so we adopt $\mathrm{Recall@FAR{\le}10^{-2}}$ as the headline biometric metric alongside
ROC-AUC, which provides a threshold-independent summary of discriminative power.
The Youden index~\cite{youden1950}, $J = \mathrm{TPR} - \mathrm{FPR}$, is used for
optimal threshold selection.
